\section{How large is selection?}
\label{sec:results}

\begin{figure}[t]
\centering
\includegraphics[width=\textwidth]{fig1-mechanism.pdf}
\caption{The mechanism in three steps. \textbf{(a)} Completed family size is
right-skewed with a coefficient of variation of \FamilySizeCV{}; the model stands
three effective types in for that distribution, at the conditional means of its
equal-probability thirds, rescaled to reproduce the dispersion exactly.
\textbf{(b)} World composition drifts because higher-propensity women contribute
more of the next cohort: the high type rises from \CompositionStartShare{} of the
world population at \PhaseFiveBaseYear{} to \CompositionHighShare{} at
\ProjectionEndYear{}. \textbf{(c)} The resulting selection multiplier, observed
fertility divided by the environment path alone, with the band showing
\MechanismDraws{} draws from the stated ranges of all
\MechanismParameterCount{} mechanism parameters.}
\label{fig:mechanism}
\end{figure}

\subsection{The size of the effect, step by step}

The cleanest benchmark is not the UN's own path, because that path already
contains an assumed fertility recovery whose justification is exactly what is in
dispute. Instead the benchmark holds that today's fertility already embodies
today's conditions: the decline already visible in the data continues to its
observed floor and then stops, with no assumed rebound and no assumed further
deterioration. Selection is then switched on against that background, and the
components are added one at a time (\Cref{fig:ladder}).

\begin{figure}[t]
\centering
\includegraphics[width=0.92\textwidth]{fig2-ladder.pdf}
\caption{World population in \ProjectionEndYear{} as each component is added.
Grey is the benchmark environment with no selection. Orange steps are the
mechanism: measured mainstream selection first, then the two named
high-fertility groups. Red is the illustrative stress test, a further
\PressureStressTest{} per decade decline in the fertility environment after
\PressureFromYear{}, which is a scenario knob and not a forecast. The
\LadderPressureLoss{}-billion size of that last step, against
\LadderMainstreamGain{} billion for the mechanism itself, is why this paper's
primary result is the boundary in \Cref{sec:boundary} rather than any single
number on this chart.}
\label{fig:ladder}
\end{figure}

Measured mainstream selection --- the two independently sourced parameters, with
no named groups and no assumed change in conditions --- raises world population
in \ProjectionEndYear{} from \LadderBenchmark{} billion to \LadderMainstream{}
billion. That is \LadderMainstreamGain{} billion, or
\LadderMainstreamGainPercent{}, from a mechanism
that is entirely absent from conventional projections. Observed world fertility
ends \RaceSelectionEffectPercent{} above the environment path that produced it,
and rises above it detectably from about \SelectionOvertakesYear{}
(\Cref{fig:mechanism}c).

The mechanism is a level shift, not a return to growth. Under the benchmark
environment with mainstream selection the world still peaks and then runs close
to flat; selection changes where the curve settles, not its shape. This is the
distinction \citet{arenberg2022} press against models built from heritability
alone, and it survives here because the projection carries absolute
age-specific fertility rather than a transmission coefficient by itself.

\subsection{Named groups are a small part of the story}

Adding the two named high-fertility groups --- with their measured fertility,
their measured retention, and one scenario knob routing their leavers --- moves
the world result by \LadderNamedGroupsGain{} billion. The world composition at
\ProjectionEndYear{} is \CompositionGroupShare{} named-group.

This deserves emphasis because it inverts the usual telling. Public discussion of
selection on fertility is almost entirely about visible high-fertility
minorities, and this model treats them generously: it gives them their measured
fertility, it lets high retention compound for four generations, and its
assumption that each group starts with its host country's age structure
understates their growth, so the shares reported here are lower bounds. They
still contribute \NamedGroupShareOfMechanism{} of what unlabelled variation inside
ordinary populations contributes. The reason is arithmetic rather than sociological. A
group at a tenth of a per cent of a national population has four generations to
compound before it is large enough to matter to a world total, whereas the
mainstream's dispersion acts on everybody at once from the first year.

Where a group is already large it does matter, and locally it dominates. In
Israel, whose Haredi population is \HarediShare{} of the country at the base
year, the named group reaches \IsraelGroupShare{} of the population by
\ProjectionEndYear{} and the country's selection multiplier reaches
\SelectionEffectIsrael{} --- half again as much fertility as the environment path
alone would give. In the United States the Amish reach
\UnitedStatesGroupShare{}, contributing to a national multiplier of
\SelectionEffectUnitedStates{}. Neither of these is a world-scale effect, and
neither should be reported as one.

\subsection{Where selection bites, and where it does not}

The same run produces very different multipliers by country
(\Cref{tab:countries}), and the ordering is informative. Selection is strongest
where fertility is already low, because the model's response depends on the
\emph{relative} dispersion of family size rather than its absolute level: a
country at 0.8 children per woman has as much proportional room for
recomposition as one at 4, and it passes through more generations in the same
number of years.

\begin{table}[t]
\centering
\small
\begin{tabular}{lrl}
\toprule
Country & Selection multiplier at \ProjectionEndYear{} & Composition of the effect \\
\midrule
Israel          & \SelectionEffectIsrael{}        & mainstream plus a large named group \\
United States   & \SelectionEffectUnitedStates{}  & mainstream plus a small named group \\
Nigeria         & \SelectionEffectNigeria{}       & mainstream only \\
Japan           & \SelectionEffectJapan{}         & mainstream only \\
Germany         & \SelectionEffectGermany{}       & mainstream only \\
Korea           & \SelectionEffectKorea{}         & mainstream only \\
\bottomrule
\end{tabular}
\caption{Observed fertility divided by the environment path alone, at
\ProjectionEndYear{}, in the run that includes both selection and the
\PressureStressTest{} stress test. The multiplier is a statement about
composition and is unaffected by the level of the environment path, for the
reason given in \Cref{sec:separability}.}
\label{tab:countries}
\end{table}

\subsection{The mechanism's own uncertainty}

Drawing all \MechanismParameterCount{} mechanism parameters from their stated
ranges \MechanismDraws{} times gives a selection multiplier at
\ProjectionEndYear{} between \RaceEnsembleEffectLow{} and
\RaceEnsembleEffectHigh{} at 90\%, and a world population in the stress-test
scenario between \RaceEnsembleLow{} and \RaceEnsembleHigh{} billion. Selection
overtakes the environment in \RaceEnsembleOvertakeShare{} of those draws.

The direction is therefore robust and the magnitude is not, and it is worth being
precise about which parameter is responsible. Almost all of the width comes from
the empirical dispersion of family size: holding everything else fixed, the
endpoints of the measured range, \FamilySizeCVLow{} and \FamilySizeCVHigh{}, move
the stress-test world result from \CVLowWorld{} to \CVHighWorld{} billion. That is a genuine
empirical uncertainty in a quantity that is measurable, unlike the environment
path, and narrowing it is tractable work: it needs completed-cohort parity
distributions from more than the small set of countries that publish them.
